% =====================================================================
% DRAFT — insert into 05-responsibility-architecture.tex
% Placement: after \subsubsection{Two Chatbots, Two Merits Results}
%            before \subsection{Corporate and Individual Responsibility Run in Parallel}
% Body length: ~1,450 words (~4 printed pages with footnotes)
% =====================================================================

\subsection{The Strongest Case Against Enterprise-Facing Recourse}

The front-door rule attracts three serious objections, and they arrive from
different directions. The first is that the architecture is unnecessary
because ordinary tort law can already reach AI harms. The second is that
expanding enterprise exposure raises the total cost of compensating injuries
and leaves the intended beneficiaries worse off. The third is that many AI
harms should not be actionable at all, so a rule guaranteeing a defendant
guarantees the wrong thing. Each objection is strong, and each narrows the
proposal rather than defeating it. The narrowing is stated below.

\subsubsection{Sufficiency: Conventional Doctrine May Already Reach These Harms}

Kenneth Abraham and Catherine Sharkey argue that tort law is more capable here
than the skeptics allow. On their account the black-box problem does not defeat
liability, because the governing standards are outcome dependent: they turn on
the reasonableness of the defendant's conduct rather than on an actor's state
of mind or on reconstruction of the model's internal operation. Objective
reasonableness, the products defective-design standard, and intentional torts
such as invasion of privacy and defamation can each be applied to
first-generation AI cases involving vehicles and chatbots with conventional or
minimally adjusted principles.\lrfootnote{See \artcite{abrahamSharkey2027}; see
also \artcite{ramakrishnanTortFrontier2026} (surveying the tort-theoretic stakes
of frontier AI harms).} If that is right, a statutory front door is redundant
machinery.

Much of this Article agrees. This Part opened by conceding that the obstacle is
not an absence of doctrine, and nothing proposed here supplies a new liability
standard. The disagreement concerns stage rather than substance. Abraham and
Sharkey answer which test governs once the defendant and the relevant conduct
are identified. The architecture proposed here addresses how a claimant reaches
that point.

The gap becomes visible in the deployment structures mapped in Part IV. An
objective standard of reasonableness still requires evidence of what the
defendant did. Where one enterprise develops, configures, distributes, and
operates a system, that evidence has a single custodian and the conventional
account works without supplementation. Where control is divided, it does not.
The claimant observes an output. To plead conduct, the claimant needs the
configuration record. To obtain the configuration record, the claimant must
name a party who holds it. To name that party, the claimant must already know
which of several enterprises selected the model, wrote the system instructions,
supplied the tools, set the objective, or declined to stop operation. The
standard is available; the facts that feed it are not. The production rule
breaks that circle without altering the standard.

The objection nevertheless earns a limit, and the limit improves the proposal.
The two components of the front door do different work and should be evaluated
separately. In a vertically integrated deployment, the identification component
is close to declaratory of existing attribution law and adds cost without adding
access; \emph{Garcia} required no statute to name Character Technologies. It is
the production component that operates there, and the identification component
that operates in the divided architectures where an injured person can
presently be defeated by the org chart. A legislature could sensibly enact the
second obligation broadly and the first only for deployments in which
development, configuration, and operation are separated.

\subsubsection{Cost: Expanded Liability Raises the Price of Compensation}

Mark Geistfeld identifies a cost that liability scholarship routinely omits.
Mainstream theory favors strict enterprise liability on the ground that
internalizing injury costs induces reasonably safe practices while liability
insurance compensates victims. Geistfeld's objection is that a commercial
defendant's liability insurance routinely indemnifies losses the plaintiff has
already insured under first-party coverage such as health insurance. Expanding
liability multiplies that duplication, and right-holders ultimately pay for it
through prices. A safety gain purchased at a larger increase in the cost of
insuring against injury leaves the protected class worse off. His conclusion is
that an effective regime combines negligence liability, narrowly tailored strict
liability, and administrative regulation rather than a general
expansion.\lrfootnote{See \artcite{geistfeldInsurability2026}.}

That conclusion is closer to the present proposal than to its opposite. This
Article does not propose strict enterprise liability, and says so at each step:
the claimant retains the burden of proving breach or defect, causation, legally
cognizable harm, and every element the governing law requires, and every defense
remains available. The duplication mechanism operates on the volume of
compensated claims. A rule that leaves the substantive elements untouched does
not enlarge that volume. It changes who must be named and what must be produced.

Two costs survive that answer and should be acknowledged. First, the production
obligation is incurred in cases the enterprise ultimately wins, and that cost is
correlated across deployments in a way ordinary litigation exposure is not.
Second, model-level failures may themselves be correlated---one version behaving
the same way across a million deployments---which is precisely the condition
under which pooling fails. For that class of risk, insurance may be the wrong
instrument, and a bounded compensation fund, a channeling rule, or a statutory
cap may serve better than either liability or coverage. That is an institutional
design question of the kind Part II identified, and it is a legislature's to
answer.

The \Charter{} also cuts the other way on insurability. Insurers price what they
can observe. A deployment with named risk owners, tested rollback, logged
incidents, and independent review is underwritable; an undocumented deployment
of the same system is not. Aviation is the working example: a mandatory
organization-wide safety management system coexists with a functioning insurance
market in part because the safety architecture generates the loss data
underwriting requires.\lrfootnote{See \statcite{faaSMSPart5}; \bluecite{faaSMS2026}.}
Documentation duties are ordinarily counted as a cost of a liability regime.
Here they are also a precondition of the actuarial tractability whose absence
Geistfeld correctly identifies as the deeper problem.

\subsubsection{Scope: Some AI Harms Should Have No Remedy}

Zachary Henderson presses the opposite concern. He argues that categorically
applying strict liability to AI harms would be a mistake, that AI activities are
not monolithic, and that many such harms should have no tort remedy at
all---a position he grounds in economic efficiency, corrective justice, and civil
recourse alike. On his account, insurance markets can fill compensatory voids for
non-tortious harm, and domains presenting catastrophic potential call for
tailored, domain-specific legislation rather than a general
rule.\lrfootnote{See \artcite{hendersonRobotLiability2026}.}

The architecture is orthogonal to that claim, and saying so costs nothing.
Whether a loss is actionable is a substantive judgment about the cause of action.
Whether an injured person can identify the enterprise and obtain the deployment
record is a question of attribution and proof. A legislature persuaded by
Henderson can enact both: it can narrow the class of compensable AI harms and
still require that, within the narrowed class, the deploying enterprise face the
claimant and produce the record. Henderson's own preferred instrument for
high-risk domains---tailored, domain-specific legislation---is the instrument this
Part proposes, and the ``high-impact deployment'' designation is an attempt to
draw the domain rather than to abolish it.

One element of the objection does reach the proposal directly. A named defendant
and a production obligation generate settlement pressure independent of the
merits, because discovery cost is itself a bargaining asset. The proposal's own
trigger supplies the first answer: the obligation attaches only after a claimant
pleads a recognized injury and a plausible material nexus to a covered
deployment, and proportionality continues to govern
scope.\lrfootnote{\statcite{frcp26}.} A legislature could add more---an early
vetting step, staged production, or fee-shifting for claims that fail at the
threshold. What it should not do is treat the absence of an identifiable
defendant as a screening device. Screening by information asymmetry filters on
the claimant's investigative resources rather than on the merits of the claim.

\subsubsection{What the Objections Change}

The three objections leave the architecture standing in a narrower form than
this Part first stated. The identification component is close to declaratory
where a deployment is integrated and does its real work where control is
divided. The production component is the operative reform; it should be
triggered by a pleading threshold, bounded by proportionality, and paired with
confidentiality protection. The substantive law of recovery is left wherever the
enacting legislature wants it, including narrower than present law. And the
charter is defended not only as an accountability instrument but as one of the
few available means of making a correlated and poorly quantified risk observable
enough to price. None of the three objections requires the machine to become a
defendant. Each is an argument about which existing persons should answer, on
what showing, and at what cost---which is the question this Article says courts
were always asking.
